% !TEX program = xelatex
% Build:  latexmk -xelatex main.tex      (จัดการ toc + bibtex ให้อัตโนมัติ)
%
% แบบฟอร์ม: การเสนอเค้าโครงโครงงานรายวิชา
% สาขาวิชาวิทยาการคอมพิวเตอร์ วิทยาลัยการคอมพิวเตอร์ มหาวิทยาลัยขอนแก่น

\documentclass[12pt,a4paper]{article}

% ---------------------------------------------------------------
% Fonts / language  (XeLaTeX เท่านั้น)
% ---------------------------------------------------------------
\usepackage{fontspec}
% ไฟล์ต้นแบบ (เค้าโครงโครงงานdatasci.docx) ใช้ TH SarabunPSK 16pt
% เครื่องนี้ไม่มี PSK มีแต่ TH Sarabun New ซึ่งเป็นตัวเดียวกันคนละรุ่น (metric ตรงกัน)
% 12pt x 1.333 = 16pt
\setmainfont{TH Sarabun New}[Scale=1.333]

% ตัดบรรทัดภาษาไทย -- ถ้าไม่ใส่ ย่อหน้าภาษาไทยจะทะลุขอบกระดาษ
\XeTeXlinebreaklocale "th"
\XeTeXlinebreakskip = 0pt plus 1pt

% ---------------------------------------------------------------
% Page layout
% ---------------------------------------------------------------
% ขอบกระดาษตามไฟล์ต้นแบบ: บน/ซ้าย 1.5 นิ้ว, ล่าง/ขวา 1 นิ้ว
\usepackage[a4paper,left=1.5in,right=1in,top=1.5in,bottom=1in]{geometry}
\usepackage{setspace}
\singlespacing              % ต้นแบบใช้ระยะบรรทัดเดี่ยว
\linespread{1.15}           % เผื่อที่ให้สระบน + วรรณยุกต์ไม่ชนบรรทัดบน
\setlength{\parindent}{1.27cm}
\setlength{\parskip}{0pt}

% ---------------------------------------------------------------
% Common content packages
% ---------------------------------------------------------------
\usepackage{graphicx}
\graphicspath{{figures/}}
\usepackage{array}
\renewcommand{\arraystretch}{1.3}   % เผื่อความสูงบรรทัดในตารางให้สระ/วรรณยุกต์
\usepackage{tabularx}
\usepackage{multirow}
\usepackage{booktabs}
\usepackage{amsmath,amssymb}
\usepackage{caption}
\usepackage{float}
\usepackage{enumitem}
\usepackage{url}

% ---------------------------------------------------------------
% หัวข้อ: ขนาดเท่าตัวอักษรปกติ ตัวหนา มีจุดหลังเลขข้อ ("1. ชื่อหัวข้อโครงงาน")
% ตามไฟล์ต้นแบบ -- ไม่ใช้หัวข้อตัวใหญ่แบบ default ของ LaTeX
% ---------------------------------------------------------------
\usepackage{titlesec}
\titleformat{\section}{\normalsize\bfseries}{\thesection.}{0.5em}{}
\titleformat{\subsection}{\normalsize\bfseries}{\thesubsection}{0.5em}{}
\titlespacing*{\section}{0pt}{1.2em}{0.4em}
\titlespacing*{\subsection}{1.27cm}{0.6em}{0.3em}

% ---------------------------------------------------------------
% หัวข้อย่อยแบบ 3.1 / 5.1 / 6.1 ตามแบบฟอร์ม
% ใช้:  \begin{subitems} \item ... \end{subitems}
% ---------------------------------------------------------------
\newlist{subitems}{enumerate}{1}
\setlist[subitems]{label=\thesection.\arabic*,labelindent=1.27cm,leftmargin=*,
                   itemsep=0.2em,topsep=0.2em,parsep=0pt}

% ---------------------------------------------------------------
% การอ้างอิง (APA, backend = BibTeX เพราะเครื่องนี้ไม่มี biber)
% ---------------------------------------------------------------
% numberedbib = ให้หัวข้อบรรณานุกรมมีเลขข้อ (กลายเป็น "11. เอกสารอ้างอิง")
\usepackage[numberedbib]{apacite}
\bibliographystyle{apacite}
% apacite ตั้ง \refname ใหม่ตอน \begin{document} จึงต้อง override ด้วย \AtBeginDocument
\AtBeginDocument{\renewcommand{\refname}{เอกสารอ้างอิง}}

% ---------------------------------------------------------------
% Links (โหลดท้ายสุด)
% ---------------------------------------------------------------
\usepackage[hidelinks,unicode]{hyperref}

\begin{document}

% ส่วนหัวของแบบฟอร์ม -- กรอกชื่อและรหัสประจำตัวตรง \dotfill
\begin{center}
    \textbf{การเสนอเค้าโครงโครงงานรายวิชา}

    \vspace{0.5em}
    สาขาวิชาวิทยาการคอมพิวเตอร์ วิทยาลัยการคอมพิวเตอร์ มหาวิทยาลัยขอนแก่น
\end{center}

\vspace{1em}

\noindent
ชื่อ นาย / นางสาว \hspace{0.5em}\dotfill\ รหัสประจำตัว \dotfill

\noindent
Mr. / Miss \dotfill

\vspace{0.8em}

\noindent
ชื่อ นาย / นางสาว \hspace{0.5em}\dotfill\ รหัสประจำตัว \dotfill

\noindent
Mr. / Miss \dotfill

\vspace{1.5em}


\section{ชื่อหัวข้อโครงงาน}

\noindent
ภาษาไทย \hspace{0.5em} การศึกษาความสัมพันธ์ระหว่างพฤติกรรมการใช้โซเชียลมีเดีย\\
\hspace*{4.2em} กับสุขภาวะทางจิตของนักศึกษา

\vspace{0.5em}

\noindent
ภาษาอังกฤษ \hspace{0.5em} A Study of the Relationship between Social Media\\
\hspace*{5.6em} Usage Behavior and Mental Well-being among\\
\hspace*{5.6em} University Students

\section{หลักการและเหตุผล}

% บอกถึงปัญหา และความเป็นมาของโครงงาน แนวความคิดเบื้องต้นของโครงงาน
% เหตุผลและความจำเป็น

\section{วัตถุประสงค์ของโครงงาน}

% ระบุให้ชัดเจนว่าทำโครงงานเพื่ออะไร เช่น จะสร้างอะไร ใช้แก้ปัญหาอะไร

\begin{subitems}
    \item เพื่อ \dotfill
    \item เพื่อ \dotfill
    \item เพื่อ \dotfill
\end{subitems}

\section{ทฤษฎีและผลงานวิจัยที่เกี่ยวข้อง}

% การศึกษาทฤษฎี รวมทั้งงานวิจัยที่เกี่ยวข้องกับโครงงานที่จะทำ
% อ้างอิงด้วย \cite{key} โดย key มาจาก references.bib

\section{วิธีดำเนินการวิจัย}

% เขียนขั้นตอนการทำโครงงาน เช่น การศึกษาค้นคว้าทฤษฎี การศึกษาเครื่องมือที่จะใช้
% การกำหนดขอบเขตและเป้าหมายของงานวิจัย การจัดเตรียมข้อมูล การออกแบบระบบ
% การพัฒนาระบบ การทดสอบ สรุปผล เขียนรายงาน และการเสนอผลงาน

\begin{subitems}
    \item \dotfill
    \item \dotfill
\end{subitems}

\section{ขอบเขตและข้อจำกัดของการวิจัย}

% ขอบเขต คือ สิ่งที่สามารถทำได้ และทำได้ดีที่สุด
% ข้อจำกัด คือ สิ่งที่ไม่สามารถทำได้
% เช่น สิ่งที่จะสร้างคืออะไร มีความสามารถทำอะไรได้บ้าง ใช้งานกับระบบใด ใช้อะไรในการพัฒนา

\begin{subitems}
    \item \dotfill
    \item \dotfill
\end{subitems}

\input{chapter/07-location.tex}
\section{ประโยชน์ที่คาดว่าจะได้รับ}

% ระบุผลลัพธ์จากงานวิจัย หรือประโยชน์ต่าง ๆ ที่เกิดจากการนำงานวิจัยไปใช้ต่อไปในอนาคต

\begin{subitems}
    \item \dotfill
    \item \dotfill
\end{subitems}

\section{แผนและระยะเวลาดำเนินการ}

% ตารางการดำเนินงานและระยะเวลา -- ทำเครื่องหมายช่วงเวลาด้วย $\bullet$ หรือระบายช่อง

\begin{table}[H]
    \centering
    \begin{tabular}{|p{6cm}|c|c|c|c|c|c|}
        \hline
        \multirow{2}{*}{\centering กิจกรรม} & \multicolumn{6}{c|}{เดือนที่} \\
        \cline{2-7}
                                            & 1 & 2 & 3 & 4 & 5 & 6 \\
        \hline
        & & & & & & \\
        \hline
        & & & & & & \\
        \hline
        & & & & & & \\
        \hline
        & & & & & & \\
        \hline
    \end{tabular}
\end{table}

\section{งบประมาณ}

\begin{table}[H]
    \centering
    \begin{tabular}{|p{9cm}|r|}
        \hline
        \multicolumn{1}{|c|}{รายการ} & \multicolumn{1}{c|}{จำนวนเงิน (บาท)} \\
        \hline
        \textbf{หมวดวัสดุอุปกรณ์} & \\
        \hspace{1em}ค่าวัสดุสำนักงาน (กระดาษ ปากกา และอื่น ๆ) & \\
        \hspace{1em}ค่าวัสดุคอมพิวเตอร์ (แผ่นดิสก์ ซีดี และอื่น ๆ) & \\
        \hline
        \textbf{หมวดค่าใช้สอย} & \\
        \hspace{1em}ค่าถ่ายเอกสาร & \\
        \hspace{1em}ค่าจัดรูปเล่ม & \\
        \hline
        \textbf{หมวดค่าใช้จ่ายอื่น ๆ} & \\
        \hspace{1em} & \\
        \hline
        \multicolumn{1}{|r|}{\textbf{รวมทั้งสิ้น}} & \\
        \hline
    \end{tabular}
\end{table}

% หัวข้อ "11. เอกสารอ้างอิง" ถูกพิมพ์โดย \bibsection ที่นิยามไว้ใน main.tex
% รายการอ้างอิงมาจาก references.bib -- ต้อง \cite{} ในเนื้อหาจึงจะขึ้น
% (หรือใช้ \nocite{*} เพื่อแสดงทุกรายการในไฟล์)
\nocite{*}
\bibliography{references}


\end{document}
